% Section 6, part 1 (parts 2-10 follow by \input). Label prefix: mod. No new macros.
\section{The model: topological recursion on the model curve}
\label{sec:mod}

Topological recursion is applied here, and only here, to the curve $x=\Wv(p)$, $y=p$ on $\bbP^1$, $v\neq 0$; nothing about the quintic enters. The Model Lemma on $\sK$ (\cref{thm:mod-MLK}) follows from the identification with topological recursion (\cref{thm:mod-tr}), a leg dictionary (\cref{prop:mod-LD,prop:mod-dii}), a lattice theorem (\cref{thm:mod-lattice}), L-regularity (\cref{thm:mod-Lreg}), the Gap Theorem (\cref{thm:mod-gap}) and Theorem~V (\cref{thm:mod-V}). The inputs (Q) are those of \cref{sec:int-inputs}, each cited where used.

\subsection{Conventions; the model as topological recursion}
\label{sec:mod-tr}

\begin{conv}[Topological recursion]
\label{conv:mod-kernel}
Let $x,y$ be meromorphic functions on $\bbP^1$ such that $\dd x$ has only simple zeros, the \emph{branch points} $\varsigma_b$, and $\dd y\neq 0$ there; let $\iota_b$ be the local involution of $x$ near $\varsigma_b$ ($x\circ\iota_b=x$, $\iota_b\neq\mathrm{id}$). Put $\omega_{0,1}:=y\,\dd x$ and $\omega_{0,2}(p_1,p_2):=\dd p_1\,\dd p_2/(p_1-p_2)^2$. For $2h-2+n+1>0$ the correlators $\omega_{h,n+1}$ are defined by the recursion, with $J=(p_1,\dots,p_n)$,
\begin{equation}
\label{eq:mod-kernel}
\begin{gathered}
\omega_{h,n+1}(p_0,J)=\sum_b\Res_{q\to\varsigma_b}K_b(p_0,q)\Big[\omega_{h-1,n+2}(q,\iota_bq,J)+{\sum}'\,\omega_{h_1,|I_1|+1}(q,I_1)\,\omega_{h_2,|I_2|+1}(\iota_bq,I_2)\Big],\\
K_b(p_0,q):=\frac{1}{2}\,\frac{\int_{\iota_bq}^{q}\omega_{0,2}(p_0,\cdot)}{\big(y(q)-y(\iota_bq)\big)\,\dd x(q)},
\end{gathered}
\end{equation}
where ${\sum}'$ runs over $h_1+h_2=h$ and $I_1\sqcup I_2=J$ with the two terms containing $\omega_{0,1}$ omitted. Here $q$ is a point of the curve. For $h\geq 2$ the free energy is defined by $(2-2h)\FTR_h:=\sum_b\Res_{\varsigma_b}\Phi_{y\dd x}\,\omega_{h,1}$, where $\dd\Phi_{y\dd x}=y\,\dd x$. This is the convention of \cite[(2-26)--(2-29)]{Dunin-Barkowski:2015caa} and of \cite[(2.1)]{Hock25}. The kernel printed in \cite[(4-3), (4-9)]{Eynard:2007kz} equals the kernel of \cite[(1.1)]{ABDKS25} with $\omega^{(0)}_{1,0}=-y\,\dd x$, and both are $-K_b$. Since $\omega_{h,n}$ is built with $2h-2+n$ kernels, the correlators $W^{(h)}_n$ of \cite{Eynard:2007kz} and $\omega^{(h)}_{n,0}$ of \cite{ABDKS25} are
$\omega_{h,n}=(-1)^nW^{(h)}_n=(-1)^n\omega^{(h)}_{n,0}$ and $\FTR_h=-F^{(h)}$, with $F^{(h)}$ of \cite[Def.~4.3 (4-10)]{Eynard:2007kz}; in particular our $\omega_{0,3}$ is the negative of \cite[Thm.~4.1 (4-20)]{Eynard:2007kz}. The same dictionary is recorded in \cref{sec:sgn-tr} (\cref{tab:sgn-tr}).
\end{conv}

By \cite[Thm.~4.2, Thm.~4.6]{Eynard:2007kz}, translated by the dictionary above, every $\omega_{h,n}$ with $2h-2+n>0$ is a symmetric meromorphic multidifferential whose poles, in each variable, lie at the branch points only; in particular it has no poles on the diagonals.

\begin{defn}[The model curve]
\label{def:mod-curve}
Let $v\in\bbC$, $v\neq 0$, fix a fifth root $\varsigma$ of $v$, and let $x:=\Wv(p)=p^4/4+v/p$ and $y:=p$ on $\bbP^1$ with the coordinate $p$. Then $\dd x=(p^5-v)\,\dd p/p^2$; the branch points are $\varsigma_b:=\zeta^{-b}\varsigma$, $b\in\bbZ/5$; $x''(\varsigma_b)=5\varsigma_b^2\neq 0$; the critical values $x(\varsigma_b)=\frac54\zeta^b\varsigma^4$ are pairwise distinct; $\dd y=\dd p\neq 0$; and $p=0$, $p=\infty$ are the poles of $x$ and of $y\,\dd x$. The five points of a fibre $x^{-1}(X)$ are the roots of $T^5-4XT+4v$. The weights are $p:1$, $v:5$, $x:4$, $\bar p:=v/p:4$, $\log p:0$. The correlators $\omega_{h,n}$ and $\FTR_h$ of this curve are those of \cref{conv:mod-kernel}; their homogeneity in $v$ is \cref{lem:mod-scale}~(ii).
\end{defn}

\begin{defn}[Pairings with legs]
\label{def:mod-pair}
The \emph{leg class} is $\cL:=\bbQ[v][p]+\bbQ(v)[x]$. For $2h-2+n>0$ and legs $\phi_1,\dots,\phi_n\in\cL$ put
\[
\langle\phi_1\cdots\phi_n\rangle^{\mathrm{TR}}_h:=\sum_{b_1,\dots,b_n}\Res_{p_1=\varsigma_{b_1}}\cdots\Res_{p_n=\varsigma_{b_n}}\prod_{l=1}^n\phi_l(p_l)\,\omega_{h,n}(p_1,\dots,p_n),
\]
the \emph{residue pairing}; for one slot we write $\langle\phi,w\rangle:=\sum_b\Res_{\varsigma_b}\phi\,w$. For legs in $\bbQ[v][p]$ the residue pairing equals the pairing $(2\pi\ri)^{-1}\oint_{|p|=R}$ on a circle with $R>|\varsigma|$, since such a leg is holomorphic at $p=0$, where $\omega_{h,n}$ is holomorphic too. Let $\mathfrak r$ be a ray from $0$ to $\infty$ avoiding the branch points. The \emph{log leg} $-\log p$ is paired by
\[
\langle\cdots,-\log p,\cdots\rangle^{\mathrm{TR}}_h:=\int_{p'\in\mathfrak r}\omega_{h,n}(\dots,p',\dots),
\]
and several log legs by iterated integrals along rays; in a mixed pairing the slots with legs in $\cL$ are paired by residues and the other slots along rays. These integrals converge and do not depend on the rays: in each variable $\omega_{h,n}$ is holomorphic off the branch points (the paragraph after \cref{conv:mod-kernel}), and its residues there vanish (\cref{lem:mod-xfun} with $f=1$, whose proof does not use pairings), and moving a ray across $\varsigma_b$ changes the integral by $2\pi\ri$ times such a residue. Moreover, if $\eta$ is a meromorphic $1$-form on $\bbP^1$ with poles only at the branch points and zero residues there, then $\eta=\dd G$ with $G$ rational and holomorphic at $0$ and $\infty$, and
\[
\int_{\mathfrak r}\eta=G(\infty)-G(0)=\sum_b\Res_{\varsigma_b}(-\log p)\,\eta,
\]
for any branch of $\log p$ near each $\varsigma_b$: indeed $\Res_{\varsigma_b}(\log p\,\dd G)=-\Res_{\varsigma_b}G\,\dd p/p$, and the residue theorem for the rational form $G\,\dd p/p$, whose other poles are $0$ and $\infty$, gives $-\sum_b\Res_{\varsigma_b}G\,\dd p/p=G(0)-G(\infty)$. So a log leg may equally be paired by residues at the branch points.
\end{defn}

\begin{defn}[Dual descendant functions]
\label{def:mod-xi}
Let $D_x:=\dd/\dd x=\frac{p^2}{p^5-v}\frac{\dd}{\dd p}$ acting on functions, and for $a\in\bbZ/5$
\[
\xi_a:=\varepsilon'_a\,v^{\alpha_a}\frac{p^{m_a}}{p^5-v},\quad (m_a,\varepsilon'_a,\alpha_a)=(4,-1,0),\ (0,1,1),\ (1,1,0),\ (2,-1,0),\ (3,-1,0)\quad(a=0,\dots,4),
\]
and $\xi_{a,d}:=D_x^d\xi_a$ for $d\geq 0$. (The signs $\varepsilon'_a$ are local, not the $\varepsilon_b$ of \cref{def:set-model}.) Each $\xi_{a,d}$ is rational with poles only at the branch points, holomorphic at $p=0$, and $O(p^{-1-4d})$ at $p=\infty$.
\end{defn}

\begin{defn}[Ancestor vector of a leg]
\label{def:mod-ancvec}
For a leg $\phi\in\cL$, or $\phi=-\log p$, or a $\bbQ[v]$-linear combination of these, put $q_\phi(z):=\sum_{a,d}\langle\phi,\dd\xi_{a,d}\rangle H^az^d$; a leg $\phi$ with $q_\phi=X$ is called a \emph{TR leg} of $X$. Here $\langle-\log p,\dd\xi_{a,d}\rangle=\xi_{a,d}(\infty)-\xi_{a,d}(0)$ by \cref{def:mod-pair}.
\end{defn}

The ancestor correlators $\langle\prod_i\psi^{d_i}H^{a_i}\rangle_h(v)$ of $\Omega^{\mathrm{mod}}_v$ are those of \cref{def:set-modcohft}, extended multilinearly to vectors $q_1,\dots,q_n\in\bbH[z]$ (\cref{conv:tra-forms}).

\begin{thm}[The model as topological recursion]
\label{thm:mod-tr}
Let $v\neq 0$ and $2h-2+n>0$. The correlators of the model curve are
\begin{equation}
\label{eq:mod-tr}
\omega_{h,n}(p_1,\dots,p_n)=(-5)^{h-1}\sum_{a_1,\dots,a_n}\sum_{d_1,\dots,d_n\geq 0}\Big\langle\prod_{i=1}^n\psi^{d_i}H^{a_i}\Big\rangle_h(v)\,\prod_{i=1}^n\dd\xi_{a_i,d_i}(p_i),
\end{equation}
a finite sum, the brackets being the ancestor correlators of $\Omega^{\mathrm{mod}}_v$. Equivalently, for legs $\phi_1,\dots,\phi_n$ as in \cref{def:mod-ancvec}, $\langle q_{\phi_1}\cdots q_{\phi_n}\rangle_h=(-5)^{1-h}\langle\phi_1\cdots\phi_n\rangle^{\mathrm{TR}}_h$.
\end{thm}

% Section 6, part 2.
Throughout \cref{lem:mod-laplace,lem:mod-scale,prop:mod-Rtr}, $\kappa\neq 0$ is a rational number and $x_\kappa:=\kappa\Wv$; the curve $(\bbP^1,x_\kappa,y=p)$ has the branch points $\varsigma_j$ of \cref{def:mod-curve}. The cases used are $\kappa=4$ and $\kappa=1$. The letter $\vartheta$ denotes the Laplace variable of \cite[(2-31), (2-32)]{Dunin-Barkowski:2015caa} (called $\zeta$ there).

\begin{lem}[Formal Laplace series]
\label{lem:mod-laplace}
Let $c$ be a branch point of $x_\kappa$ and $w$ a local coordinate at $c$ with $x_\kappa=x_\kappa(c)-w^2/2$. For a $1$-form $\omega=\sum_{m\geq -M}a_mw^m\,\dd w$ meromorphic at $c$ with $\Res_c\omega=0$ put
\begin{equation}
\label{eq:mod-laplace}
\begin{gathered}
\mathsf L_c[\omega](\vartheta):=\sum_ma_m\,\mathsf G_m(\vartheta),\\
\mathsf G_{2n}:=(2n-1)!!\,\vartheta^{-n}\ (n\geq 0),\qquad \mathsf G_{-2n}:=\frac{(-1)^n\vartheta^n}{(2n-1)!!}\ (n\geq 1),\qquad \mathsf G_m:=0\ (m\text{ odd}),
\end{gathered}
\end{equation}
a formal Laurent series in $\vartheta^{-1}$ (the coefficient $a_{-1}$ vanishes). Then:
\begin{enumerate}
\item if $F$ is meromorphic at $c$ and $F\,\dd x_\kappa$ has zero residue at $c$, then $\mathsf L_c[\dd F]=-\vartheta\,\mathsf L_c[F\,\dd x_\kappa]$;
\item $w$ is determined by $x_\kappa$ up to $w\mapsto -w$, which multiplies $\mathsf L_c$ by $-1$;
\item let $k\in\bbC^\times$ and $v=k^5v'$, so that $x_{\kappa,v}(kp')=k^4x_{\kappa,v'}(p')$ and $\varsigma_j(v)=k\varsigma_j(v')$, and normalise the coordinates $w_j$ at $\varsigma_j$ by $\frac{\dd w_j}{\dd p}(\varsigma_j)=\sqrt{-5\kappa}\,\varsigma_j$ with one square root for all $v$; then $w_j(v)(kp')=k^2w_j(v')(p')$ and $\mathsf L_{\varsigma_j(v)}[p^m\dd p](\vartheta)=k^{m-1}\,\mathsf L_{\varsigma_j(v')}[p'^m\dd p'](k^4\vartheta)$;
\item if $\omega$ is holomorphic at $c$, then $\vartheta\frac{\dd}{\dd\vartheta}\mathsf L_c[\omega]=\vartheta\,\mathsf L_c\big[(x_\kappa-x_\kappa(c))\,\omega\big]+\frac12\mathsf L_c[\omega]$.
\end{enumerate}
\end{lem}

The integrals $\int\omega\,e^{(x_\kappa-x_\kappa(c))\vartheta}$ in \cite[(2-31), (2-32)]{Dunin-Barkowski:2015caa} are read as $\sqrt{2\pi/\vartheta}\,\mathsf L_c[\omega]$.

\begin{proof}
All four statements are linear, so it suffices to check monomials. (i) For $F=w^m$, $m\neq -2$ (for $F=w^{-2}$ the form $F\,\dd x_\kappa$ has a residue), $\dd F=mw^{m-1}\dd w$ and $F\,\dd x_\kappa=-w^{m+1}\dd w$, so the claim is $m\mathsf G_{m-1}=\vartheta\mathsf G_{m+1}$: both sides are $(2n+1)!!\,\vartheta^{-n}$ for $m=2n+1\geq 1$, $(-1)^n\vartheta^{n+1}/(2n-1)!!$ for $m=-(2n+1)$, and $0$ for even $m$. (ii) $w^2=2(x_\kappa(c)-x_\kappa)$, and $w\mapsto -w$ multiplies $w^m\dd w$ by $(-1)^{m+1}$, while only even $m$ contribute. (iii) $w_j(v)(kp')$ and $k^2w_j(v')(p')$ are coordinates at $\varsigma_j(v')$ with the same square $2k^4(x_{\kappa,v'}(\varsigma_j(v'))-x_{\kappa,v'}(p'))$ and the same derivative $k^2\sqrt{-5\kappa}\,\varsigma_j(v')$ there, so they agree. If $p'^m\dd p'=\sum_nb_nw'^n\dd w'$ with $w'=w_j(v')$, then $p^m\dd p=\sum_nk^{m-1-2n}b_nw^n\dd w$ with $w=k^2w'$, and $\mathsf G_{2s}(\vartheta)=k^{4s}\mathsf G_{2s}(k^4\vartheta)$, $\mathsf G_{-2s}(\vartheta)=k^{-4s}\mathsf G_{-2s}(k^4\vartheta)$ give the claim. (iv) On $w^{2n}\dd w$ the left side is $-n(2n-1)!!\,\vartheta^{-n}$; as $(x_\kappa-x_\kappa(c))w^{2n}\dd w=-\frac12w^{2n+2}\dd w$, the right side is $-\frac12(2n+1)!!\,\vartheta^{-n}+\frac12(2n-1)!!\,\vartheta^{-n}$, the same; odd powers give $0$ on both sides.
\end{proof}

\begin{lem}[Scaling of the curve]
\label{lem:mod-scale}
\begin{enumerate}
\item For constants $\kappa,\mu\neq 0$ let $\omega^{(\kappa,\mu)}_{h,n}$ and $F^{(\kappa,\mu)}_h$ be the correlators and free energies of $(\bbP^1,\kappa\Wv,\mu p,\omega_{0,2})$. Then $\omega^{(\kappa,\mu)}_{h,n}=(\kappa\mu)^{2-2h-n}\omega_{h,n}$ for $2h-2+n>0$, and $F^{(\kappa,\mu)}_h=(\kappa\mu)^{2-2h}\FTR_h$ for $h\geq 2$.
\item For $k\in\bbC^\times$ let $m_k(p):=kp$. Then $m_k^*\,\omega^{(k^5v)}_{h,n}=k^{5(2-2h-n)}\,\omega^{(v)}_{h,n}$ and $\FTR_h(k^5v)=k^{5(2-2h)}\FTR_h(v)$, where the superscript is the parameter of the curve. Consequently a pairing of $\omega_{h,n}$ with the legs $p^{m_1},\dots,p^{m_j}$ in $j$ slots and with $-\log p$ in the other $n-j$ slots is a function $C(v)$ with $C(k^5v)=k^{5e}C(v)$, $e=2-2h-n+\sum_lm_l/5$.
\end{enumerate}
\end{lem}

\begin{proof}
(i) The branch points, the involutions $\iota_b$ and $\omega_{0,2}$ are unchanged, and the kernel becomes $K_b/(\kappa\mu)$. By induction on $2h-2+n$, starting from the exponent $0=2-0-2$ of $\omega_{0,2}$: every term of \eqref{eq:mod-kernel} for $\omega_{h,n+1}$ is one kernel times $\omega_{h-1,n+2}$, of exponent $2-2h-n$, or times $\omega_{h_1,|I_1|+1}\omega_{h_2,|I_2|+1}$, of exponent $(1-2h_1-|I_1|)+(1-2h_2-|I_2|)=2-2h-n$, so the total exponent is $2-2h-(n+1)$. For the free energy, $\Phi_{y\dd x}$ becomes $\kappa\mu\,\Phi_{y\dd x}$. (ii) The automorphism $m_k$ of $\bbP^1$ satisfies $m_k^*W_{k^5v}=k^4\Wv$, $m_k^*p=kp$ and $m_k^*\omega_{0,2}=\omega_{0,2}$, and maps the branch points of $\Wv$ to those of $W_{k^5v}$. The recursion is natural under isomorphisms of curves, so $m_k^*\omega^{(k^5v)}_{h,n}$ is the correlator of $(\bbP^1,k^4\Wv,kp,\omega_{0,2})$, which is $k^{5(2-2h-n)}\omega^{(v)}_{h,n}$ by (i); likewise for $\FTR_h$. For the pairings, pull the residues back along $m_k$: a leg $p_l^{m_l}$ contributes $k^{m_l}$, and a ray goes to a ray, along which the pairing does not depend on the choice (\cref{def:mod-pair}).
\end{proof}

\begin{prop}[The $R$-matrix of the curve]
\label{prop:mod-Rtr}
Let $v\neq 0$ and take at the branch points of $x_\kappa$ the coordinates $w_j$ of \cite[\S2.6]{Dunin-Barkowski:2015caa}, $x_\kappa=x_\kappa(\varsigma_j)-w_j^2/2$, normalised by $\frac{\dd w_j}{\dd p}(\varsigma_j)=\sqrt{-5\kappa}\,\varsigma_j$ with one square root for all $j$. Let $e^{\mathrm{TR}}_j$, $\Delta_j$ and $\xi^j$ be the corresponding data (2-30), (2-34) of \cite[\S2.6]{Dunin-Barkowski:2015caa}: $\Delta_j^{-1/2}=\frac{\dd y}{\dd w_j}(\varsigma_j)$ and $\xi^j(q)=\int_\infty^q\omega_{0,2}(w_j,\cdot)/\dd w_j|_{w_j=0}$; and let $R_{\mathrm{TR}}$ be the matrix series defined by (2-31) read through \cref{lem:mod-laplace}, $R_{\mathrm{TR}}^{-1}(\vartheta^{-1})^j{}_i:=-\vartheta^{-1}\mathsf L_{\varsigma_j}[\dd\xi^i]$ (upper index the row). Then, as formal series in $\vartheta^{-1}$:
\begin{enumerate}
\item $R_{\mathrm{TR}}^{-1}(\vartheta^{-1})^j{}_i=\Mo(t)_{ij}$ with $t=-1/(\kappa\varsigma^4\vartheta)$, for all $i,j$; consequently
\begin{equation}
\label{eq:mod-Rtr}
R_{\mathrm{TR}}(\psi)=\Mo\big(\psi/(\kappa\varsigma^4)\big)
\end{equation}
in the basis $e^{\mathrm{TR}}_j$, and $R_{\mathrm{TR}}$ is symplectic in this orthonormal basis;
\item (2-32) holds: $\sum_kR_{\mathrm{TR}}^{-1}(\vartheta^{-1})^i{}_k\Delta_k^{-1/2}=\mathsf L_{\varsigma_i}[\dd y]$, both sides being $\varsigma^{-1}\Rhat^{\mathrm m,(i)}_0(t)/\sqrt{-5\kappa}$. In the frozen basis $\bar e_j:=e^{\mathrm{TR}}_j/\sqrt{-5\kappa}$ the unit $\sum_k\Delta_k^{-1/2}e^{\mathrm{TR}}_k$ is $\hat 1:=\sum_k\bar e_k/\varsigma_k$, and the dilaton leaves of the graph formula carry the translation $\psi\big(\hat 1-\Mo^{-1}(\psi/(\kappa\varsigma^4))\hat 1\big)=\psi\,\hat 1-\psi R_{\mathrm{TR}}^{-1}(\psi)\hat 1$ of \cite[Def.~2.13]{PPZ15}.
\end{enumerate}
\end{prop}

The sign of the dilaton leaves is that of \cite[\S2.3]{DuninBarkowski:2012bw} and of \cite[Prop.~2.12, Def.~2.13]{PPZ15}; \cite[\S2.3]{Dunin-Barkowski:2015caa} prints the opposite sign. With legs $R^{-1}(\psi)$ the sign is not a convention: it is the one for which the action preserves the unit (\cref{sec:sgn-R}).

% Section 6, part 3.
\begin{proof}[Proof of \cref{prop:mod-Rtr}]
Write $\sqrt{-5\kappa}$ for the root of the statement and $\mathsf L_j:=\mathsf L_{\varsigma_j}$; indices run over $\bbZ/5$.

\emph{Step 1 (integration by parts).} By (2-34), $\dd\xi^i(q)=\frac{\dd p}{\dd w_i}(\varsigma_i)\,\dd q/(\varsigma_i-q)^2$, so $\xi^i(q)=-1/\big(\sqrt{-5\kappa}\,\varsigma_i(q-\varsigma_i)\big)$. It has at most a simple pole at $\varsigma_j$, so $\xi^i\dd x_\kappa$ is holomorphic there, and \cref{lem:mod-laplace}~(i) gives $R_{\mathrm{TR}}^{-1}(\vartheta^{-1})^j{}_i=\mathsf L_j[\xi^i\dd x_\kappa]$.

\emph{Step 2 (partial fractions).} Since $\dd x_\kappa=\kappa(p^5-v)\,\dd p/p^2$ and $(p^5-v)/(p-\varsigma_i)=\sum_{k=0}^4\varsigma_i^kp^{4-k}$, with the local letters $J^{(j)}_m:=\mathsf L_j[p^m\dd p]$, $\tilde J^{(j)}_m:=\sqrt{-5\kappa}\,J^{(j)}_m$:
\begin{equation}
\label{eq:mod-Rtr-pf}
R_{\mathrm{TR}}^{-1}(\vartheta^{-1})^j{}_i=\frac15\sum_{k=0}^4\varsigma_i^{k-1}\,\tilde J^{(j)}_{2-k}.
\end{equation}

\emph{Step 3 (the five series are model rows).} Fix $j$, write $J_m:=J^{(j)}_m$, $z:=-1/(\kappa\vartheta)$ (so $\vartheta x_\kappa=-\Wv/z$) and $t:=z/\varsigma^4$. \emph{(Mod2).} \Cref{lem:mod-laplace}~(i) with $F=p^m$ gives $mzJ_{m-1}-J_{m+3}+vJ_{m-2}=0$; in particular $J_3=vJ_{-2}$. \emph{(Mod1).} By \cref{lem:mod-laplace}~(iii) with $k=\varsigma$, $v'=1$, $J_m(v;\vartheta)=\varsigma^{m-1}J_m(1;\varsigma^4\vartheta)$, so $5v\,\partial_vJ_m=(m-1)J_m+4\vartheta\partial_\vartheta J_m$. By \cref{lem:mod-laplace}~(iv), $\vartheta\partial_\vartheta J_m=\vartheta\kappa\big(\tfrac14J_{m+4}+vJ_{m-1}\big)-\tfrac54\vartheta\kappa\varsigma_j^4J_m+\tfrac12J_m$; eliminating $J_{m+4}=(m+1)zJ_m+vJ_{m-1}$ by (Mod2), using $\kappa\vartheta=-1/z$ and inserting the result into (Mod1) gives $z\partial_vJ_m=-J_{m-1}+J_m/\varsigma_j$ (as $\varsigma_j^5=v$). Put $\tilde Y_m:=e^{-\Wv(\varsigma_j)/z}J_m$ (a formal symbol, with the Leibniz rule); as $\partial_v[\Wv(\varsigma_j(v))]=1/\varsigma_j$, $z\partial_v\tilde Y_m=-\tilde Y_{m-1}$. Hence $Y_b:=\varepsilon_b\tilde Y_{\nu_b}$ satisfy $z\partial_vY_b=b^{\mathrm m}_bY_{b+1}$ (for $b=1$ use $Y_2=-\tilde Y_3=-v\tilde Y_{-2}$ by (Mod2)), the derivation of \cref{prop:set-modrows}~(a) for these series. Write $Y_b=e^{U_j/z}\varsigma^{\nu_b-1}F_b$ with $U_j:=-\Wv(\varsigma_j)$ and $F_b:=\varepsilon_b\varsigma^{1-\nu_b}J_{\nu_b}$. By the homogeneity above, $F_b=\varepsilon_bJ_{\nu_b}(1;\varsigma^4\vartheta)$ is a power series in $t$, and the computation of \cref{prop:set-modrows}~(b) shows that $(F_b)_b$ solves the $T$-system of colour $j$. \emph{Initial values.} $J_m=\varsigma_j^{m-1}/\sqrt{-5\kappa}+O(\vartheta^{-1})$ and $(\varsigma_j/\varsigma)^{\nu_b-1}=\zeta^{j(b+1)}$, so $\sqrt{-5\kappa}F_b(0)=\varepsilon_b\zeta^{j(b+1)}=\Rhat^{\mathrm m,(j)}_b(0)$ (\cref{def:set-M0}). \emph{Uniqueness.} By \cref{lem:gau-colmod}~(i), (ii), $\Rhat^{\mathrm m,(j)}=\zeta^jF^{(j)}$ also solves the colour-$j$ system, and by \cref{prop:set-modrows}~(b) (pivot $4n\zeta^{4j}$) a power-series solution is determined by the constant term of its component $b=0$. Hence $\tilde J^{(j)}_{m_b}=\varepsilon_b\varsigma^{m_b-1}\Rhat^{\mathrm m,(j)}_b(t)$ for $(m_0,\dots,m_4)=(0,-1,-2,2,1)$, $m_b\equiv -b\pmod 5$ (for $b=2$, $\tilde J_{-2}=\tilde J_3/v$).

\emph{Step 4 (assembly).} As $b$ runs over $\bbZ/5$, $k:=2-m_b$ runs over $0,\dots,4$, and $\varsigma_i^{k-1}\varsigma^{m_b-1}=\zeta^{-i(b+1)}$. By Step 3 and \eqref{eq:mod-Rtr-pf}, $R_{\mathrm{TR}}^{-1}(\vartheta^{-1})^j{}_i=\frac15\sum_b\varepsilon_b\zeta^{-i(b+1)}\,\Rhat^{\mathrm m,(j)}_b(t)=\sum_b(\Psi_0^{-1})_{ib}\Rhat^{\mathrm m,(j)}_b(t)=\Mo(t)_{ij}$, since $(\Psi_0)_{bi}=\varepsilon_b\zeta^{i(b+1)}$ has the inverse $\frac15\varepsilon_b\zeta^{-i(b+1)}$.

\emph{Step 5 (inverse).} With the upper index as row, Step 4 reads $R_{\mathrm{TR}}^{-1}(\psi)=\Mo\big(-\psi/(\kappa\varsigma^4)\big)^T$, $\psi=\vartheta^{-1}$. By \cref{lem:set-M0sympl}, $\Mo(-t)^T=\Mo(t)^{-1}$, which gives \eqref{eq:mod-Rtr} and symplecticity. This proves (i).

(ii) By (2-30), $\Delta_k^{-1/2}=1/(\sqrt{-5\kappa}\varsigma_k)$. By (i), $1/\varsigma_k=\zeta^k/\varsigma$ and $\sum_k\zeta^{-kb}=5\delta_{b,0}$, the left side of (2-32) is $\sum_k\Mo(t)_{ki}/(\sqrt{-5\kappa}\,\varsigma_k)=\Rhat^{\mathrm m,(i)}_0(t)/(\sqrt{-5\kappa}\,\varsigma)$, and the right side is $J^{(i)}_0=\Rhat^{\mathrm m,(i)}_0(t)/(\sqrt{-5\kappa}\varsigma)$ by Step 3. The unit is $\sum_k\Delta_k^{-1/2}e^{\mathrm{TR}}_k=\sum_k\bar e_k/\varsigma_k$. Given (2-32), in the graph formula of \cite[Thm.~4.1]{DuninBarkowski:2012bw} a dilaton leaf is a marked point of its vertex carrying $\psi^k$, $k\geq 2$, with the weight $\coef{z^{k-1}}\bigl(-R(-z)\bigr)$ applied to the unit \cite[\S2.3]{DuninBarkowski:2012bw}, in total $\psi\bigl(1-R(-\psi)\bigr)$ applied to the unit; by \cref{lem:set-pushins}~(ii) at every vertex these points give the vertex integrals of the translation by this series in the push-forward form of \cite[Def.~2.5]{PPZ15} and \cite[\S2.3]{Dunin-Barkowski:2015caa}; and $R(-z)$ there, which also decorates the ordinary leaves, is $R^{-1}(z)$ of \cite[\S2.3]{Dunin-Barkowski:2015caa}. Finally $R_{\mathrm{TR}}^{-1}(\psi)=\Mo(\psi/(\kappa\varsigma^4))^{-1}$ by (i).
\end{proof}

\begin{proof}[Proof of \cref{thm:mod-tr}]
\emph{Step 1 (the quoted theorem).} \cite[Thm.~4.1]{DuninBarkowski:2012bw}, as recalled in \cite[\S2.6, (2-30)--(2-35)]{Dunin-Barkowski:2015caa}, with the action of \cite[\S2.3]{Dunin-Barkowski:2015caa} (legs $R^{-1}(\psi)$, edges as in \cite[Def.~2.2]{PPZ15}, dilaton leaves with the translation of \cite[Def.~2.13]{PPZ15}, in the push-forward form of \cite[Def.~2.5]{PPZ15} by the proof of \cref{prop:mod-Rtr}~(ii)), gives for a compact curve with simple branch points, $\dd y\neq 0$ there, and the Bergman kernel: if the matrix $R$ of (2-31) is symplectic in the orthonormal basis and (2-32) holds, then for $2h-2+n>0$
\begin{equation}
\label{eq:mod-235}
\omega_{h,n}=\sum_{i_1,\dots,i_n}\sum_{d_1,\dots,d_n\geq 0}\int_{\overline{\mathcal M}_{h,n}}\alpha_{h,n}(e_{i_1},\dots,e_{i_n})\prod_{l=1}^n\psi_l^{d_l}\;\prod_{l=1}^n\dd\big((\dd/\dd x)^{d_l}\xi^{i_l}\big)(p_l),
\end{equation}
where $\alpha=R.\TFT^{\mathrm{TR}}$ and $\TFT^{\mathrm{TR}}_{h,n}(e_j,\dots,e_j)=\Delta_j^{(2h-2+n)/2}$, zero on mixed arguments. We apply it to $f:=4\Wv$, the superpotential \cite[(8-1)]{Dunin-Barkowski:2015caa} with $n=4$, $a=(0,0,0,0,4v)$: the critical points $p^5=v$ are simple with distinct critical values, $\dd y\neq 0$, and the shift $p_{\mathrm{Dub}}=p$ is trivial ($a_1=0$); the compatibility condition \cite[Prop.~4.3 (1)]{Dunin-Barkowski:2015caa} holds, since the curve is rational and $f$ has a pole at $\infty$ ($\dd y/\dd f=p^2/(4(p^5-v))$ has poles only at the zeros of $\dd f$); and by the sentence preceding \cite[Thm.~8.1]{Dunin-Barkowski:2015caa} this curve satisfies (2-30)--(2-33). The hypotheses on $R$ are \cref{prop:mod-Rtr} at $\kappa=4$, which gives $R^{4}(\psi)=\Mo(\psi/(4\varsigma^4))$, $\Delta^{4}_j=-20\varsigma_j^2$, $\xi^{(4,j)}=-1/(\sqrt{-20}\,\varsigma_j(p-\varsigma_j))$; so \eqref{eq:mod-235} holds for $(\bbP^1,4\Wv,p)$ with $\alpha^{4}=R^{4}.\TFT^{4}$.

\emph{Step 2 (transfer to $\kappa=1$).} Take $\sqrt{-20}=2\sqrt{-5}$ and identify $e^4_j\leftrightarrow e^1_j$. Then $\xi^{(4,j)}=\xi^{(1,j)}/2$, $\dd/\dd x_4=\frac14\dd/\dd x_1$, $\TFT^4_{h,n}=2^{2h-2+n}\TFT^1_{h,n}$, the unit of $\TFT^4$ is half that of $\TFT^1$, and $R^4(\psi)=R^1(\psi/4)$ with $R^1(\psi):=\Mo(\psi/\varsigma^4)$. In the graph sum a vertex with $n_v$ special points gives $2^{2h_v-2+n_v}$, each dilaton leaf $2^{-1}$, and $\sum_v(2h_v-2+n_v)=2h-2+n+\#\{\text{dilaton leaves}\}$, so the TFT contributes $2^{2h-2+n}$. The substitution $R\mapsto R(\cdot/c)$ multiplies every leg, dilaton and edge term by $c^{-D}$, $D$ its cohomological degree on $\overline{\mathcal M}_{h,n}$ (for edges, $E_{R(\cdot/c)}(\psi',\psi'')=c^{-1}E_R(\psi'/c,\psi''/c)$), so $[R(\cdot/c).\TFT]_{h,n}=\sum_Dc^{-D}[R.\TFT]^{(D)}_{h,n}$. In \eqref{eq:mod-235} only $D=3h-3+n-\sum_ld_l$ contributes, and leg $l$ gives $4^{-d_l}2^{-1}$ times the leg for $\kappa=1$. Hence $\omega^{(4,1)}_{h,n}=2^{2h-2+n}4^{-(3h-3+n)}2^{-n}\,\mathrm{RHS}^1=4^{2-2h-n}\,\mathrm{RHS}^1$, and $\omega^{(4,1)}_{h,n}=4^{2-2h-n}\omega_{h,n}$ by \cref{lem:mod-scale}~(i). So \eqref{eq:mod-235} holds on $x=\Wv$ with $\alpha^1=\Mo(\cdot/\varsigma^4).\TFT^1$, $\Delta^1_j=-5\varsigma_j^2$, $\xi^{(1,j)}=-1/(\sqrt{-5}\,\varsigma_j(p-\varsigma_j))$.

\emph{Step 3 (the frozen basis and the factor $(-5)^{h-1}$).} With $\bar e_j:=e^1_j/\sqrt{-5}$ and $\tilde\xi^j:=\sqrt{-5}\,\xi^{(1,j)}=-1/(\varsigma_j(p-\varsigma_j))$ the summand of \eqref{eq:mod-235} is unchanged, $\TFT^1_{h,n}(\bar e_j,\dots,\bar e_j)=(-5)^{h-1}\varsigma_j^{2h-2+n}$, $\eta(\bar e_j,\bar e_k)=-\delta_{jk}/5$, and the unit is $\hat 1$ (\cref{prop:mod-Rtr}~(ii)). Identifying $\bar e_j$ with the normalised canonical vector $\tilde e_j$ of $\TFT^{\mathrm m}_v$ ($\hat\Delta_j^{1/2}=\varsigma_j$, \cref{def:set-modcohft}), $\TFT^1_{h,n}=(-5)^{h-1}\TFT^{\mathrm m}_{v,h,n}$, with metric $-\frac15$ times that of $\TFT^{\mathrm m}_v$ and the same unit. For a CohFT $\Omega$ with metric $\eta$ and $\lambda\neq 0$, $\lambda^{h-1}\Omega_{h,n}$ is a CohFT for $\eta/\lambda$ with the same unit, and $R.(\lambda^{h-1}\Omega)=\lambda^{h-1}(R.\Omega)$: each vertex gives $\lambda^{h_v-1}$, each edge $\lambda$, and $\sum_v(h_v-1)+\#\{\text{edges}\}=h-1$. With $\lambda=-5$: $\alpha^1=(-5)^{h-1}\big(\Mo(\cdot/\varsigma^4).\TFT^{\mathrm m}_v\big)$ on the frozen basis.

\emph{Step 4 (the flat basis).} By \cref{def:set-modcohft}, $\Omega^{\mathrm{mod}}_v=(\Psi^{\mathrm m}_v)_*(\Mo(z/\varsigma^4).\TFT^{\mathrm m}_v)$, with $\eta(H^b,\tilde e_i)=\varepsilon_b\varsigma^{\nu_b-1}\zeta^{i(b+1)}$ (\cref{def:set-M0}, \cref{prop:set-modrows}~(b)). As $H^{3-a}/5$ is $\eta$-dual to $H^a$, $\sum_i\tilde e_i\otimes\tilde\xi^i=\sum_aH^a\otimes\xi'_a$ with $\xi'_a:=\frac15\sum_i\eta(H^{3-a},\tilde e_i)\tilde\xi^i$. Using $\nu_{3-a}=2-\nu_a$ and $\zeta^{-ia}=(\varsigma_i/\varsigma)^a$,
\[
\xi'_a=-\frac{\varepsilon_{3-a}}{5}\,\varsigma^{-\nu_a-a}\sum_i\frac{\varsigma_i^a}{p-\varsigma_i},\qquad\sum_i\frac{\varsigma_i^a}{p-\varsigma_i}=\begin{cases}5p^4/(p^5-v),&a=0,\\ 5vp^{a-1}/(p^5-v),&1\leq a\leq 4,\end{cases}
\]
by partial fractions. With $\nu_a+a=0,0,5,5,5$ and $\varepsilon_{3-a}=1,-1,-1,1,1$ this gives $\xi'_a=\xi_a$ of \cref{def:mod-xi}, and \eqref{eq:mod-235} becomes \eqref{eq:mod-tr}; the sum is finite since a correlator vanishes unless $\sum_ld_l\leq 3h-3+n$. Pairing each slot with a leg (\cref{def:mod-pair}) gives the second form.
\end{proof}

% Section 6, part 4.
\begin{lem}[Functions of $x$ pair to zero]
\label{lem:mod-xfun}
Let $2h-2+n>0$, $n\geq 1$, let $p_2,\dots,p_n$ be points of $\bbP^1$ that are not branch points, and let $f$ be holomorphic near the five critical values $x(\varsigma_b)$ (for instance $f\in\bbQ(v)[X]$, $v\neq 0$ fixed). Then for each $b$,
\[
\Res_{p=\varsigma_b}f(x(p))\,\omega_{h,n}(p,p_2,\dots,p_n)=0 .
\]
Consequently every pairing of $\omega_{h,n}$ in which one slot is paired with $f(x)$ by the residue pairing, the other slots being paired by residues at the branch points, on circles or along rays, vanishes; legs $\phi,\phi'\in\cL$ with $\phi-\phi'\in\bbQ(v)[x]$ have the same residue pairings; and (with $f=1$) $\omega_{h,n}$ has zero residues at the branch points in each variable.
\end{lem}

\begin{proof}
\emph{The trace identity.} For $X\in\bbP^1$ let $p^1(X),\dots,p^5(X)$ be the roots of $T^5-4XT+4v$, the points of $x^{-1}(X)$. We claim $\mathrm{Tr}\,\omega(X):=\sum_{i=1}^5\omega_{h,n}(p^i(X),p_2,\dots,p_n)=0$. For $q'\notin x^{-1}(X)$, $\prod_i(q'-p^i(X))=4q'\,(x(q')-X)$, so $\sum_i\omega_{0,2}(p^i(X),q')=\dd X\,\dd x(q')/(X-x(q'))^2$, and hence $\sum_i\int_{\iota_bq}^q\omega_{0,2}(p^i(X),\cdot)=\dd X\big[(X-x(q))^{-1}-(X-x(\iota_bq))^{-1}\big]=0$. Take disjoint $\iota_b$-stable small discs $U_b\ni\varsigma_b$ avoiding $p_2,\dots,p_n$, $U:=\bigcup_bU_b$. For $p_1\notin U$, \eqref{eq:mod-kernel} writes $\omega_{h,n}(p_1,J)$ as $\sum_b(2\pi\ri)^{-1}\oint_{\partial U_b}K_b(p_1,q)$ times the bracket, which does not depend on $p_1$ (the path in $K_b$ stays in $U_b$). For $X\notin x(\overline U)$, summing over $i$ under the integrals gives $\mathrm{Tr}\,\omega(X)=0$, hence $\mathrm{Tr}\,\omega=0$ by meromorphy (as in \cite[App.~A, (A-1), (A-5)]{Eynard:2007kz}).

\emph{The residues.} Put $X_b:=x(\varsigma_b)$. The fibre $x^{-1}(X_b)$ is $\varsigma_b$ (local degree $2$) and three points $q_1,q_2,q_3$ of local degree $1$, none a branch point (the critical values are distinct) and none $0$ or $\infty$. So $\eta:=f(x)\,\omega_{h,n}(\cdot,J)$ is holomorphic at the $q_k$ (\cref{conv:mod-kernel}). The preimage of a small circle around $X_b$ is a union of loops around the points of the fibre, so $\Res_{X_b}\mathrm{Tr}(\eta)=\sum_{q\in x^{-1}(X_b)}\Res_q\eta$, and $\mathrm{Tr}(\eta)=f(X)\,\mathrm{Tr}\,\omega=0$ gives $\Res_{\varsigma_b}\eta=0$.

\emph{Consequences.} By symmetry the slot carrying $f(x)$ may be the first; the contours of the other slots can be taken disjoint from small discs around the branch points, and there are no diagonal poles, so the residue in the first slot may be taken first, and it vanishes. (Cf.\ \cite[Thm.~4.4 (4-23), Cor.~4.1 (4-26)--(4-29), App.~A (A-25)]{Eynard:2007kz}, not used here.)
\end{proof}

\subsection{The leg dictionary and the residue identity}
\label{sec:mod-legs}

\begin{defn}[The legs]
\label{def:mod-legs}
Let $\phi^b_0:=p,\ \log p,\ v/p,\ p^3,\ p^2$ for $b=0,\dots,4$. For $F$ in the span of the $p^j(\log p)^l$ ($j\in\bbZ$, $l\geq 0$) with coefficients in $\bbQ[v]$, let $I(F)$ be the termwise antiderivative of $-F\,\dd x=-F(p^3-vp^{-2})\,\dd p$ without a term $p^0(\log p)^0$ (that of $p^j(\log p)^l\dd p$ is $p^{j+1}$ times a polynomial in $\log p$ if $j\neq -1$, and $(\log p)^{l+1}/(l+1)$ if $j=-1$). Put $\phi^b_{k+1}:=I(\phi^b_k)$. Then $D_x\phi^b_{k+1}=-\phi^b_k$, and $\phi^b_k$ is homogeneous of weight $\wt(\phi^b_k)=w_b+4k$, where $w_b:=\nu_b+1$, that is $(w_0,\dots,w_4)=(1,0,4,3,2)$.
\end{defn}

\begin{prop}[The leg dictionary]
\label{prop:mod-LD}
Let $\varrho=(\varrho_0,\dots,\varrho_4):=(-1,1,-1,-3,-2)$. For all $a,b\in\bbZ/5$ and $k,d\geq 0$,
\begin{equation}
\label{eq:mod-LD}
\langle\phi^b_k,\dd\xi_{a,d}\rangle=\sum_c\Res_{c}\phi^b_k\,\dd\xi_{a,d}=-\varrho_b\,\delta_{ab}\,\delta_{dk},
\end{equation}
$c$ running over the branch points and $\log p$ being any branch near each of them. Hence $q_{\phi^b_k}=-\varrho_b\,z^kH^b$, and the TR leg of $v^mz^kH^b$ (the insertion $v^m\psi^kH^b$, \cref{conv:tra-forms}) is $-\varrho_b^{-1}v^m\phi^b_k$. In particular the unit $H^0$ has the TR leg $y=p$; $\Phi_{y\dd x}=p^5/5-v\log p=-\phi^0_1$ has $q_{\Phi_{y\dd x}}=-z\,H^0$; $-\log p$ has $q_{-\log p}=H$; the Euler-dilaton vector $\se=z\,H^0+vH$ has the TR leg $-p^5/5$; and $q_{f(x)}=0$ for $f\in\bbQ(v)[x]$.
\end{prop}

\begin{proof}
\emph{Step 1 (branches and functions of $x$).} Each leg is analytic near every $c\neq 0$. A change of branch $\log p\mapsto\log p+2\pi\ri N$ changes $\phi^b_k$ by a polynomial $\delta^b_k$ in $x$ with constant coefficients, by induction: $\delta^b_0\in\{0,2\pi\ri N\}$, and $\dd\phi^b_{k+1}=-\phi^b_k\,\dd x$ survives the substitution, so $\delta^b_{k+1}+\int_0^x\delta^b_k$ is constant. For $f$ holomorphic near $c$, repeated integration by parts gives $\Res_cf(x)\,\dd\xi_{a,d}=\pm\Res_cf^{(d+1)}(x)\,\xi_a\dd x=0$, as $\xi_a\dd x=\varepsilon'_av^{\alpha_a}p^{m_a-2}\dd p$ is holomorphic at $c$. So the pairings do not depend on the branches, and $q_{f(x)}=0$.
\emph{Step 2 (lowering).} For $d\geq 1$, $\xi_{a,d}\dd\phi=(D_x\phi)\,\dd\xi_{a,d-1}$, so integration by parts gives $\langle\phi,\dd\xi_{a,d}\rangle=-\langle D_x\phi,\dd\xi_{a,d-1}\rangle$, that is $\langle\phi^b_{k+1},\dd\xi_{a,d}\rangle=\langle\phi^b_k,\dd\xi_{a,d-1}\rangle$.
\emph{Step 3 ($d=0$, $k\geq 1$).} $\langle\phi^b_k,\dd\xi_a\rangle=\sum_c\Res_c\phi^b_{k-1}\,\xi_a\dd x=0$, as $\xi_a\dd x$ is holomorphic at $c$.
\emph{Step 4 ($k=0$, $d\geq 1$).} $\langle\phi^b_0,\dd\xi_{a,d}\rangle=-\sum_c\Res_cD_x^d(\xi_a)\,\dd\phi^b_0$. This rational form is holomorphic at $p=0$ ($D_x$ maps functions holomorphic at $0$ to $O(p^2)$, and $\dd\phi^b_0=O(p^{-2})\dd p$) and $O(p^{1-4d})\dd p$ at $\infty$; by the residue theorem the sum vanishes.
\emph{Step 5 ($k=d=0$).} Write $\dd\phi^b_0=\gamma_bp^{n_b}\dd p$ with $(\gamma_b,n_b)=(1,0),(1,-1),(-v,-2),(3,2),(2,1)$ (local constants). Since $\Res_cp^N\dd p/(p^5-v)=c^{N-4}/5$ and $\sum_cc^N=5v^{N/5}$ if $5\mid N$ and $0$ otherwise, $\langle\phi^b_0,\dd\xi_a\rangle=-\varepsilon'_a\gamma_bv^{\alpha_a}\,\frac15\sum_cc^{m_a+n_b-4}$, non-zero only if $m_a+n_b\equiv 4\pmod 5$, that is $a=b$; for $a=b=0,\dots,4$ the values are $1,-1,1,3,2=-\varrho_b$.
\emph{Step 6.} If $k=d$, Steps 2 and 5 give $-\varrho_b\delta_{ab}$; if $k>d$, Steps 2 and 3 give $0$; if $k<d$, Steps 2 and 4 give $0$. This is \eqref{eq:mod-LD}. The consequences: $\phi^0_0=p$, $-\varrho_0=1$; $\phi^0_1=-(p^5/5-v\log p)$, so $q_{\Phi_{y\dd x}}=\varrho_0z\,H^0=-zH^0$; $q_{\log p}=-\varrho_1H=-H$; and the TR leg of $\se$ is $\phi^0_1-v\phi^1_0=-p^5/5$.
\end{proof}

\begin{cor}[Free energies of the model and of the curve]
\label{cor:mod-trF}
For $h\geq 2$ let $\Fm_h(v)=\langle\ \rangle_h(v)=\int_{\overline{\mathcal M}_h}\Omega^{\mathrm{mod}}_{v;h,0}$ be the genus-$h$ free energy of the model CohFT (\cref{def:set-modcohft}), the value at the base point $-z1$ of the dilaton-shifted potential $\overline F_h$ (\cref{def:set-modcohft,def:set-graphsum}). Then $\Fm_h=(-5)^{1-h}\FTR_h$.
\end{cor}

\begin{proof}
By \cref{conv:mod-kernel}, $(2-2h)\FTR_h=\langle\Phi_{y\dd x}\rangle^{\mathrm{TR}}_h$, for any branch of $\log p$ (\cref{lem:mod-xfun}); only this definition is used (cf.\ \cite[Def.~4.3 (4-10), Thm.~4.7 (4-30)]{Eynard:2007kz}). By \cref{thm:mod-tr} and \cref{prop:mod-LD}, $\langle\Phi_{y\dd x}\rangle^{\mathrm{TR}}_h=-(-5)^{h-1}\langle\psi H^0\rangle_h$, and the dilaton equation for the flat unit $H^0$ (\cref{lem:set-dilaton}~(i), (ii)) gives $\langle\psi H^0\rangle_h=(2h-2)\Fm_h$. Comparing, $(2-2h)\FTR_h=(2-2h)(-5)^{h-1}\Fm_h$.
\end{proof}

\input{sections/s06_model_5}



